\documentclass[12pt, a4paper]{article}
\usepackage{ctex}
\usepackage{amsmath, amssymb}
\usepackage{geometry}
\usepackage{listings}
\usepackage{xcolor}
\usepackage{booktabs}
\usepackage{graphicx}
\usepackage{hyperref}
\usepackage{float}
\usepackage{enumitem}
\geometry{margin=2.5cm}

\definecolor{codebg}{RGB}{245,245,245}
\lstset{basicstyle=\ttfamily\small, breaklines=true, frame=single, backgroundcolor=\color{codebg}}

\title{\textbf{多媒体技术与应用课程作业报告}\\[4pt]\large 个人主页搭建与博客设计}
\author{刘浩东 \quad PB23010347 \\ 中国科学技术大学}
\date{\today}

\begin{document}
\maketitle

\section{个人主页介绍}

\subsection{主页主要内容}

本个人主页定位为学术技术方向，展示以下核心内容：

\begin{itemize}[leftmargin=*]
    \item \textbf{个人介绍}：姓名、学校、个人简介，位于页面 Hero 区域的左侧卡片中。
    \item \textbf{技术兴趣标签}：Python、机器学习、C/C++ 等，以标签徽章形式展示。
    \item \textbf{统计数据}：文章数、项目数、最近更新时间等动态统计，位于 Hero 区域下方。
    \item \textbf{开源项目展示}：4 个开源项目卡片（SVD 图片压缩、昆虫分类、Seam Carving、华东杯数学建模 C 题），包含描述与技术栈标签，支持展开/收起。
    \item \textbf{博客入口}：文章列表展示，点击进入独立阅读页面。
    \item \textbf{关于我}：学校、所在地、兴趣方向等卡片式介绍。
\end{itemize}

主页支持\textbf{搜索功能}，可在导航栏右侧搜索文章关键词与项目名称，下拉显示匹配结果。

\subsection{页面结构与功能说明}

页面整体采用垂直单页布局，导航栏固定置顶：

\begin{table}[H]
\centering
\caption{页面结构说明}
\begin{tabular}{p{3cm}p{9cm}}
\toprule
\textbf{区块} & \textbf{功能} \\
\midrule
导航栏 & 固定顶部，包含 Logo、分区导航链接、搜索框，移动端折叠为汉堡菜单 \\
Hero 区域 & 左侧个人卡片（头像 + 姓名 + 介绍 + 标签 + 社交链接 + 统计数据），右侧看板娘图片 \\
关于我 & 文字介绍 + 四张信息卡片（学校 / 所在地 / 兴趣方向 / 分享内容）\\
项目 & 四张项目卡片，默认展示 3 个，支持"更多项目"展开/收起 \\
文章 & 文章列表，默认展示 5 篇，支持"更多文章"展开，点击进入独立阅读页面 \\
文章独立页 & 左侧目录导航（自动提取 h2/h3），中间文章正文 + KaTeX 公式渲染 + 评论 \\
页脚 & 版权信息 \& 社交链接，返回顶部按钮 \\
\bottomrule
\end{tabular}
\end{table}

\subsection{主页访问链接}
\begin{itemize}
    \item \textbf{线上}：\href{https://lanqilhd.github.io}{https://lanqilhd.github.io}
    \item \textbf{源码}：\href{https://github.com/lanqilhd/lanqilhd.github.io}{https://github.com/lanqilhd/lanqilhd.github.io}
\end{itemize}

\section{博客内容介绍}

\subsection{博客主题}

第一篇博客主题为《从一条线到高维空间：彻底搞懂支持向量机（SVM）》，从零开始系统性介绍 SVM 的数学原理。

\subsection{主要内容概述}

文章内容的完整逻辑链如下：

\begin{enumerate}[leftmargin=*]
    \item \textbf{问题引入}：通过邮件分类的直观例子，引出二分类问题与最优决策边界。
    \item \textbf{硬间隔 SVM}：从几何直觉到数学规划，推导 $\min_{\mathbf{w},b} \frac{1}{2}\|\mathbf{w}\|^2$ 的凸二次规划问题。
    \item \textbf{软间隔}：引入松弛变量 $\xi_i$ 和惩罚参数 $C$，解释噪声处理与泛化平衡。
    \item \textbf{升维与核技巧}：通过 RBF 核的泰勒级数展开，展示无穷维特征映射的数学本质。
    \item \textbf{对偶问题}：构造拉格朗日函数，推导对偶形式并揭示支持向量的稀疏性。
    \item \textbf{SMO 算法详解}：两变量子问题的解析求解公式、上下界裁剪、阈值更新、启发式选择策略。
    \item \textbf{多分类与总结}：OvR/OvO 策略，SVM 的实际应用场景与局限性。
\end{enumerate}

文章由笔者根据 LaTeX 代码原创编写，包含 4 张配图（线性可分示意、二维例子、核技巧可视化、模型对比）。

\subsection{博客访问方式}
\begin{itemize}
    \item \textbf{主页文章列表}：首页"文章"区域点击博客卡片
    \item \textbf{直接链接}：\href{https://lanqilhd.github.io\#/post/svm-intro}{https://lanqilhd.github.io\#/post/svm-intro}
    \item \textbf{搜索}：在导航栏搜索框输入"SVM"即可找到并跳转
\end{itemize}

\section{实现过程}

\subsection{使用的主要工具或技术}

\begin{table}[H]
\centering
\caption{技术栈清单}
\begin{tabular}{ll}
\toprule
\textbf{类别} & \textbf{技术选型} \\
\midrule
前端语言 & HTML5 + CSS3 + JavaScript (ES6) \\
数学公式渲染 & KaTeX v0.16.9 (CDN) \\
Markdown 解析 & Marked.js (CDN) \\
字体 & Google Fonts (Inter + Noto Sans SC + JetBrains Mono) \\
评论系统 & Giscus (基于 GitHub Discussions) \\
版本控制 & Git + GitHub \\
部署平台 & GitHub Pages \\
\bottomrule
\end{tabular}
\end{table}

\subsection{网页搭建和部署过程}

\begin{enumerate}[leftmargin=*]
    \item \textbf{项目初始化}：在本地创建 \texttt{index.html}、\texttt{css/style.css}、\texttt{js/main.js}、\texttt{js/blog.js} 等基础文件结构，使用无框架纯静态方案。
    \item \textbf{全局样式设计}：通过 CSS 变量定义配色、字体、间距、圆角等设计 Token，支持自动暗色模式。
    \item \textbf{Hero 区域}：采用 flex 布局，左侧为个人卡片 + 统计区（flex-column），右侧为看板娘图片，使用 \texttt{align-items: stretch} 实现左右等高。
    \item \textbf{项目区与文章区}：项目卡片使用 CSS Grid 自适配列宽；文章列表使用 flex 布局，点击通过 hash 路由（\texttt{\#/post/id}）跳转到独立阅读页。
    \item \textbf{KaTeX 集成}：在 \texttt{index.html} 加载 KaTeX CDN，在 \texttt{blog.js} 中调用 \texttt{renderMathInElement} 渲染 \$...\$ 和 \$\$...\$\$ 定界符。
    \item \textbf{搜索功能}：监听 \texttt{input} 事件，遍历文章 \texttt{posts} 数组和项目 card 的 \texttt{textContent} 进行关键词匹配，防抖 200ms。
    \item \textbf{文章独立页}：动态生成 TOC 目录导航，使用 \texttt{IntersectionObserver} 实现滚动时章节高亮。
    \item \textbf{响应式设计}：设置 4 个断点（$\ge1400$px / $\le1024$px / $\le768$px / $\le480$px），匹配超大屏、平板横屏、平板竖屏、手机四类设备。
    \item \textbf{部署上线}：推送至 GitHub 仓库 \texttt{lanqilhd/lanqilhd.github.io}，GitHub Pages 自动构建部署。
\end{enumerate}

\subsection{遇到的问题及解决方法}

\begin{table}[H]
\centering
\caption{问题与解决方案}
\begin{tabular}{p{5cm}p{7cm}}
\toprule
\textbf{问题} & \textbf{解决方案} \\
\midrule
看板娘图片与左侧区域不等高 & 将 \texttt{align-items: center} 改为 \texttt{align-items: stretch}，图片设 \texttt{height: 100\%} \\
KaTeX 数学公式渲染失败 & 将渲染调用放在 \texttt{marked.parse()} 之后，使用 \texttt{requestAnimationFrame} 延迟确保 DOM 就绪 \\
文章日期显示"-1 天" & 原因是 UTC 时间与本地时间偏差，改用 \texttt{new Date(year, month, date)} 日期级比较 \\
搜索下拉框被遮挡 & 设置 \texttt{z-index: 200}，并监听全局点击关闭下拉 \\
大文件 CSV 超过 GitHub 限制 & 接受警告（不影响推送），后续可考虑 Git LFS \\
移动端导航过于拥挤 & 设计汉堡菜单组件，CSS \texttt{transform} 动画实现展开/收起 \\
\bottomrule
\end{tabular}
\end{table}

\section{AI Agent 使用说明}

\subsection{使用了哪些 AI 工具}

\begin{itemize}
    \item \textbf{Trae IDE 内置 AI Agent}：作为主要的代码助手，运行于 Trae IDE 集成开发环境中。
\end{itemize}

\subsection{AI 辅助完成了哪些内容}

\begin{table}[H]
\centering
\caption{AI 辅助内容详表}
\begin{tabular}{p{3.5cm}p{8.5cm}}
\toprule
\textbf{工作项} & \textbf{AI 贡献说明} \\
\midrule
页面整体架构 & AI 生成初始 HTML 骨架和 CSS 变量体系 \\
Hero 区域布局 & AI 完成 flex 布局代码、响应式断点适配 \\
博客系统 & AI 编写 \texttt{blog.js} 完整逻辑（列表渲染、文章加载、hash 路由）\\
KaTeX 集成 & AI 添加 CDN 依赖并编写渲染调用代码 \\
SVM 文章 Markdown & AI 将原始 LaTeX 代码转换为 Markdown 格式并修正公式定界符 \\
搜索功能 & AI 编写防抖搜索逻辑、下拉结果渲染 \\
TOC 章节导航 & AI 编写 \texttt{generateTOC()} 函数和 IntersectionObserver 高亮逻辑 \\
汉堡菜单 & AI 编写 CSS 动画和 JS 交互逻辑 \\
CSS 响应式 & AI 编写 4 断点完整适配代码 \\
统计区动态更新 & AI 编写 GitHub API 调用和日期智能显示逻辑 \\
项目卡片内容 & AI 根据仓库文件名生成中英文项目描述和技术栈标签 \\
\bottomrule
\end{tabular}
\end{table}

\subsection{自己对 AI 生成内容做了哪些修改}

\begin{itemize}[leftmargin=*]
    \item \textbf{布局调整}：AI 初始生成的居中布局改为参考 zyyo.net 的左侧卡片 + 右侧看板娘布局，并多次微调间距、字号。
    \item \textbf{配色与字体}：统一采用 Inter + Noto Sans SC + JetBrains Mono 字体组合，配色从 AI 建议的蓝色系调整后确认。
    \item \textbf{项目内容}：AI 根据仓库文件名推测的项目描述和标签，经过人工审核和修正。
    \item \textbf{SVM 文章}：AI 转换的 Markdown 中部分公式定界符和图表路径需要修正，人工补充了部分数学推导细节。
    \item \textbf{交互细节}：项目/文章默认数量限制、展开/收起按钮文案、搜索下拉样式等 UI 细节需要根据实际体验反复调整。
    \item \textbf{响应式断点数值}：AI 建议的响应式数值经过实际测试后多次调整（如 mobile 下搜索框宽度从 140px 调至 100px）。
    \item \textbf{统计区选择}：AI 最初添加了仓库数和标签数统计，经判断后删除仓库数（GitHub 邮箱不可见）和标签数（数据少不够有意义）。
\end{itemize}

\subsection{对 AI 辅助编程的思考}

\textbf{效率提升显著}：AI Agent 能够快速生成大量标准化的 HTML/CSS/JS 代码，尤其是在重复性工作（如响应式断点适配、多元素样式覆盖）上节省了大量时间。一个包含搜索、TOC、哈希路由等功能的博客系统，如果纯手写可能需要数小时，在 AI 辅助下缩短到分钟级。

\textbf{仍需人工把控方向}：AI 的代码质量依赖于提示词的质量。在实际开发中，笔者需要清晰地描述布局意图（如"看板娘与左侧等高"）、发现问题后精准定位（如"日期显示-1天"），并在多轮迭代中持续纠正 AI 的偏差。

\textbf{细节打磨离不开人工}：AI 生成的初版代码往往在视觉细节上不够精致，颜色偏冷/偏暖、间距过小或过大、动画不够流畅等。这些都需要人通过实际浏览感受来反复微调。

\textbf{迭代式协作模式}：本次开发采用的模式是"描述需求 → AI 生成代码 → 预览 → 发现问题 → 描述修改 → AI 修正"的循环。这种模式发挥了人类的设计感知优势与 AI 的代码生成效率优势，最终实现了 1+1>2 的效果。

\section{结果展示}

\subsection{主页截图}

主页包含 Hero 区域（个人卡片 + 看板娘 + 统计数据）、关于我、项目展示、文章列表四大板块。

\subsection{博客页面截图}

点击文章后进入独立阅读页面，左侧为章节导航栏，中间为文章正文（含 KaTeX 数学公式渲染），底部为 Giscus 评论区。

\subsection{响应式布局展示}

页面在桌面端、平板端、手机端均能正常显示，平板及以下设备导航栏自动切换为汉堡菜单。

\subsection{搜索功能展示}

在导航栏搜索框输入关键词（如"SVM"、"分类"、"数学建模"），可实时搜索文章与项目，下拉显示匹配结果。

\section{总结}

本次作业完成了个人主页的搭建，包含个人介绍、开源项目展示、博客系统等核心功能。页面采用纯 HTML/CSS/JS 技术栈，通过 GitHub Pages 部署上线，支持搜索、KaTeX 公式渲染、独立文章页、章节导航等高级功能，并实现了 4 级响应式布局适配。

在开发过程中，笔者深度使用了 AI Agent 辅助编程，在需求理解、代码生成、Bug 修复、样式微调等环节与 AI 形成高效协作，最终产出了一个功能完整、设计美观的个人主页。

\end{document}
